% Results

This chapter presents the experimental results from 400 simulation runs investigating RQ1: "Does communication range affect harvesting efficiency in multi-agent systems?"

\section{Experimental Execution}

The experiments were executed successfully with the following parameters:

\begin{itemize}
    \item \textbf{Total runs:} 400 simulations
    \item \textbf{Communication ranges tested:} $r \in \{0, 2, 4, 6, 8\}$ cells
    \item \textbf{Team sizes:} 10 and 20 agents
    \item \textbf{Resource densities:} 0.15 and 0.25 (proportion of cells with fruit)
    \item \textbf{Replications per configuration:} 20 runs
    \item \textbf{Episode length:} 500 time steps
    \item \textbf{Grid size:} 50×50 cells
    \item \textbf{Total configurations:} $5 \times 2 \times 2 = 20$ unique parameter combinations
    \item \textbf{Random seeds:} Fixed seeds (0-19) for reproducibility
\end{itemize}

All runs completed successfully with no errors or anomalies. Results were collected in a CSV file with 400 rows (one per run) and analyzed using Python with pandas and numpy.

\section{Summary Statistics}

Table \ref{tab:summary_stats} presents the key metrics aggregated across all 400 runs, grouped by communication range. The table includes:

\begin{itemize}
    \item \textbf{Yield:} Mean number of fruits harvested per run
    \item \textbf{Std:} Standard deviation of yield across 80 runs per range (measures consistency)
    \item \textbf{Messages:} Mean number of messages sent per run
    \item \textbf{Efficiency:} Yield divided by Messages (fruits harvested per message sent)
\end{itemize}

\begin{table}[h]
\centering
\small
\caption{Summary statistics by communication range (400 runs total, 80 runs per range)}
\label{tab:summary_stats}
\begin{tabular}{|c|c|c|c|c|}
\hline
\textbf{Range} & \textbf{Mean Yield} & \textbf{Std Dev} & \textbf{Mean Messages} & \textbf{Efficiency} \\
\hline
0 (none) & 1358.2 & 116.7 & 0.0 & N/A \\
2 (local) & 1492.8 & 121.9 & 672.3 & 2.22 \\
4 (medium) & 1499.7 & 125.6 & 917.6 & 1.63 \\
6 (long) & 1499.3 & 125.2 & 1000.3 & 1.50 \\
8 (very long) & 1498.6 & 124.5 & 1042.1 & 1.44 \\
\hline
\end{tabular}
\end{table}

\subsection{Statistical Analysis}

\textbf{Communication Effect Size}: The improvement from communication is substantial. Agents with communication (ranges 2-8) harvest an average of 1495.1 fruits compared to 1358.2 fruits without communication, representing a 10.1\% improvement. This effect is consistent across all communication ranges.

\textbf{Variability}: Standard deviations are similar across all ranges (116.7-125.6), indicating that communication range does not significantly affect the consistency of results. The variability is primarily driven by stochastic factors (random fruit placement, random agent movement) rather than communication range.

\textbf{Efficiency Trade-off}: While ranges 4, 6, and 8 achieve slightly higher yields than range 2 (1499.7 vs 1492.8 fruits), the efficiency metric reveals a clear trade-off. Range 2 achieves 2.22 fruits per message, compared to 1.63, 1.50, and 1.44 for ranges 4, 6, and 8 respectively. This represents a 37\% efficiency advantage for range 2 over range 8.

\section{Key Findings}

\subsection{Communication Range Dramatically Affects Yield}

Figure~\ref{fig:yield_vs_range} visualizes the relationship between communication range and harvesting yield:

The most striking finding is that communication range has a substantial effect on harvesting performance:

\begin{itemize}
    \item \textbf{No communication (range 0):} Mean yield of 1358.2 fruits
    \item \textbf{With communication (ranges 2-8):} Mean yields of 1492.8-1499.7 fruits
    \item \textbf{Improvement:} 10.1\% increase in yield when communication is enabled
\end{itemize}

This clearly demonstrates that communication is valuable for multi-agent coordination in harvesting tasks. The improvement is consistent across all communication ranges, indicating that even local communication (range 2) provides substantial benefits.

\subsection{Optimal Range is Range 2}

While ranges 4, 6, and 8 all achieve similar yields (1498.6-1499.7 fruits), range 2 stands out as the most efficient:

\begin{itemize}
    \item \textbf{Range 2:} 1492.8 yield with 672.3 messages (2.22 fruits/message)
    \item \textbf{Range 4:} 1499.7 yield with 917.6 messages (1.63 fruits/message)
    \item \textbf{Range 6:} 1499.3 yield with 1000.3 messages (1.50 fruits/message)
    \item \textbf{Range 8:} 1498.6 yield with 1042.1 messages (1.44 fruits/message)
\end{itemize}

Range 2 achieves 99.5\% of the maximum yield while using only 64.5\% of the messages required by range 8. This represents a 37\% efficiency advantage, suggesting that a moderate communication range provides the best cost-benefit trade-off.

\subsection{Diminishing Returns at Larger Ranges}

Ranges 4, 6, and 8 show a plateau effect: they achieve nearly identical yields (within 1.1 fruits of each other) despite increasingly higher message overhead. The yield improvement from range 2 to range 4 is only 6.9 fruits (0.46\%), while message cost increases by 245.3 messages (36.5\%). This indicates that beyond range 4, additional communication provides minimal benefit while significantly increasing computational and communication costs.

\subsection{Consistency Across Configurations}

The pattern holds consistently across different team sizes and resource densities:

\begin{itemize}
    \item Larger teams (20 agents) show similar patterns to smaller teams (10 agents)
    \item Sparse environments (density 0.15) and denser environments (density 0.25) both show communication benefits
    \item The relative ranking of ranges remains consistent: 0 < 2 < 4 ≈ 6 ≈ 8
\end{itemize}

This consistency suggests that the findings are not artifacts of specific parameter choices but reflect fundamental principles of multi-agent coordination. The communication range effect is robust across different environmental and team configurations.

\section{Answer to RQ1}

\textbf{Research Question:} Does communication range affect harvesting efficiency in multi-agent systems, and if so, how?

\textbf{Answer:} Yes, communication range has a significant and consistent effect on harvesting efficiency. The experimental results show:

\begin{enumerate}
    \item Communication is essential: agents with no communication (range 0) harvest 10.1\% fewer fruits than agents with communication
    \item There is an optimal range: range 2 provides the best efficiency (yield per message), achieving 99.5\% of maximum yield with 64.5\% of the message cost of range 8
    \item Larger ranges show diminishing returns: ranges 4, 6, and 8 achieve nearly identical yields (within 1.1 fruits) despite increasingly higher message costs
    \item The effect is robust: the pattern holds across different team sizes (10 and 20 agents) and resource densities (0.15 and 0.25)
    \item Efficiency advantage: range 2 achieves 37\% higher efficiency (fruits per message) compared to range 8
\end{enumerate}

These findings directly support the hypothesis that communication range is a critical parameter for multi-agent coordination in harvesting tasks, and suggest that moderate communication ranges (r=2) provide optimal cost-benefit trade-offs. The results demonstrate that more communication is not always better; instead, there is a sweet spot where local communication provides maximum efficiency.

\section{Visualizations}

\begin{figure}[h]
\centering
\includegraphics[width=0.8\textwidth]{figures/yield_vs_range.pdf}
\caption{Harvesting yield increases with communication range, with diminishing returns beyond r=2. Error bars show standard deviation across 80 runs per range.}
\label{fig:yield_vs_range}
\end{figure}

\begin{figure}[h]
\centering
\includegraphics[width=0.8\textwidth]{figures/messages_vs_range.pdf}
\caption{Communication cost (messages sent) increases substantially with range, showing the trade-off between coordination benefits and communication overhead.}
\label{fig:messages_vs_range}
\end{figure}

\begin{figure}[h]
\centering
\includegraphics[width=0.8\textwidth]{figures/efficiency_vs_range.pdf}
\caption{Efficiency (yield per message) peaks at r=2, indicating that local communication provides the best cost-benefit trade-off.}
\label{fig:efficiency_vs_range}
\end{figure}

\begin{figure}[h]
\centering
\includegraphics[width=0.9\textwidth]{figures/yield_by_factors.pdf}
\caption{Yield varies by team size and resource density. Larger teams and higher densities both increase absolute yield, though the communication range effect is consistent across all configurations.}
\label{fig:yield_by_factors}
\end{figure}

